\documentclass[11pt]{article}
\usepackage[a4paper,margin=1in]{geometry}
\usepackage{amsmath,amssymb,amsthm,booktabs,enumitem,mathtools,microtype}
\usepackage[T1]{fontenc}
\usepackage{lmodern}
\usepackage[hidelinks]{hyperref}

\newtheorem{theorem}{Theorem}[section]
\newtheorem{proposition}[theorem]{Proposition}
\newtheorem{lemma}[theorem]{Lemma}
\newtheorem{corollary}[theorem]{Corollary}
\theoremstyle{definition}
\newtheorem{definition}[theorem]{Definition}
\theoremstyle{remark}
\newtheorem{remark}[theorem]{Remark}

\newcommand{\mob}{\mu}
\newcommand{\kap}{\kappa_2}
\newcommand{\Ecal}{\mathcal E}
\newcommand{\Z}{\mathbb Z}
\newcommand{\dd}{\delta}
\newcommand{\thh}{\vartheta}

\hypersetup{
  pdftitle={Square-Clock Monotonicity and Finite-Clock Nonattainment for Phasewise Chowla-Free Memory},
  pdfauthor={RH research program},
  pdfsubject={Exact recurrence, saturation, and finite-clock gap for a phasewise Mobius factor class},
  pdfkeywords={Mobius function, finite clocks, max-plus optimization, run recurrence, nonattainment}
}

\title{Square-Clock Monotonicity and Finite-Clock Nonattainment\\
for Phasewise Chowla-Free Memory}
\author{RH research program}
\date{August 7, 2026}

\begin{document}
\maketitle

\begin{abstract}
We close the finite-clock attainment question left open by RH-379, without
enlarging its factor class.  A clock $q$ is fixed before $N\to\infty$; the
admissible objects are universally distance-two-safe phasewise lag-two
tables, and their interpolation coefficient $c_{11}(r)$ is zero at every
phase.  At the square clocks
\[
 q_y=4\prod_{i\le y}p_i^2,
 \qquad A_y=\prod_{i\le y}(p_i^2-1),
 \qquad D_y=\prod_{i\le y}(p_i^2-2),
\]
let $R_\ell^{(y)}$ count positive runs of length $\ell$ in the odd support
word, and put
\[
 \Ecal_y=\sum_{\ell\ \mathrm{even}}R_\ell^{(y)},\quad
 L_y=\sum_{\ell\ \mathrm{even}}\ell R_\ell^{(y)},\quad
 M_y=\sum_{\ell\ \mathrm{odd}}(\ell-1)R_\ell^{(y)}.
\]
Deleting the new prime-square residue from replicated old runs gives the
all-order identity
\[
 \Ecal_{y+1}=(p_{y+1}^2-2)\Ecal_y+M_y.
\]
Combining it with the locked RH-374 odd-run recurrence and the exact RH-379
square-clock formula yields an exact increment.  Since
$L_y-2\Ecal_y=2R_4^{(y)}+4R_6^{(y)}+6R_8^{(y)}\ge6$, the sequence $G(q_y)$
is strictly increasing.

We then prove a deliberately special saturation law: if $q_y\mid Q$ and
$Q$ has exactly the same prime support, the mod-$4$ and mod-$9$ zero-weight
separators split the fine addition-by-two cycles into replicated finite
paths, and $G(Q)=G(q_y)$.  For arbitrary fixed $q$, taking
$Q=\operatorname{lcm}(q,q_y)$ therefore gives
\[
 G(q)\le G(q_y)<B_\infty,
 \qquad
 B_\infty-G(q)\ge
 \frac{12}{\pi^2A_y(p_{y+1}^2-1)}.
\]
Thus the RH-379 all-clock supremum is not attained by any finite clock.
This is not a growing-clock result, a general cyclic-cover theorem, an
adaptive-capacity limit, an operator or trace construction, a zero
identification, or a statement about the Riemann hypothesis.
\end{abstract}

\noindent\textbf{Keywords:} M\"obius function; square clocks; max-plus
optimization; run deletion; finite-clock nonattainment.

\section{Frozen class and predecessor inputs}

The object optimized here is exactly the RH-379 phasewise Chowla-free
memory class \cite{RH379}.  We restate the boundary because every theorem
below depends on it.

\begin{definition}[Fixed-clock phasewise class]
Fix a finite integer $q\ge1$ before taking any natural-density limit.  At
each $r\in\Z/q\Z$, choose a table
\[
 f_r:\{-1,0,1\}^2\longrightarrow\{-1,+1\},
 \qquad
 \epsilon_n=f_{n\bmod q}\bigl(\mob(n-2),\mob(n)\bigr).
\]
The phase family is universally safe if no ternary input can produce
$\epsilon_n=\epsilon_{n+2}=+1$.  It belongs to the declared class only if
the coefficient $c_{11}(r)$ of
$\mob(n-2)\mob(n)$ in the RH-379 interpolation is zero for every phase
$r$.  Let $G(q)$ be the largest absolute fixed-clock limiting correlation
over this class.
\end{definition}

The phasewise condition is essential.  If $c_{11}(r)\ne0$ at even one
phase, the limiting expansion contains a phase-weighted shift-two M\"obius
correlation not controlled by the frozen sources.  No cancellation between
different phases is assumed here.

We use three exact predecessor results.  RH-374 introduces the square-clock
support word, proves its run formula and odd-run recurrence, proves the
persistence of a length-eight run, and defines $B_\infty$ as the limit of
its strictly increasing square-clock sequence \cite{RH374}.  RH-375
identifies the same constant as the all-clock supremum in the one-site
class \cite{RH375}.  RH-379 reduces the present memory class to the three
actions $0,J,I$, proves input reflection, and establishes
\begin{equation}
 \sup_{q<\infty}G(q)=B_\infty.
 \label{eq:rh379-supremum}
\end{equation}
It did not decide finite-clock attainment.

For reference, the exact progression weights at phase $r$ are
\[
 w_r(0)=0,\qquad w_r(J)=\dd_{q,r}-\thh_{q,r},\qquad
 w_r(I)=\dd_{q,r},
\]
where $\dd_{q,r}$ is the squarefree density and $\thh_{q,r}$ is the joint
squarefree density at displacement two.  Along an addition-by-two cycle,
the compatibility matrix is
\begin{equation}
\begin{array}{c|ccc}
 &0&J&I\\ \hline
0&1&1&1\\
J&1&1&0\\
I&1&1&0
\end{array}.
\label{eq:compatibility}
\end{equation}
In particular, action $0$ is compatible with everything in both directions.
RH-379 proves that the positive three-state optimum equals $G(q)$: input
reflection realizes the negative of every limiting score, so the absolute
optimum introduces no additional case.

\section{Square-clock runs and a deletion recurrence}

Let $3=p_1<p_2<\cdots$ be the odd primes and define
\begin{equation}
 P_y=\prod_{i\le y}p_i^2,\qquad q_y=4P_y,\qquad
 A_y=\prod_{i\le y}(p_i^2-1),\qquad
 D_y=\prod_{i\le y}(p_i^2-2).
 \label{eq:square-parameters}
\end{equation}
On $\Z/P_y\Z$, set
\begin{equation}
 w_y(k)=\mathbf1_{\{p_i^2\nmid 2k+1\ \text{for every }i\le y\}}.
 \label{eq:support-word}
\end{equation}
The multiples of $9$ cut its cyclic positive runs, so every run has length
at most eight.  Write $R_\ell^{(y)}$ for the number of runs of exact length
$\ell$.  We need four statistics:
\begin{align}
 O_y&=R_1^{(y)}+R_3^{(y)}+R_5^{(y)}+R_7^{(y)},
 &\Ecal_y&=R_2^{(y)}+R_4^{(y)}+R_6^{(y)}+R_8^{(y)},
 \label{eq:OE}\\
 L_y&=2R_2^{(y)}+4R_4^{(y)}+6R_6^{(y)}+8R_8^{(y)},
 &M_y&=2R_3^{(y)}+4R_5^{(y)}+6R_7^{(y)}.
 \label{eq:LM}
\end{align}
Thus $L_y$ counts sites inside even runs, while $M_y$ is the sum of
$\ell-1$ over odd runs.  These quantities must not be conflated with one
another.

RH-374 gives, with
\[
 e_m^{(y)}=\prod_{i\le y}\left(1-\frac{m}{p_i^2}\right),
\]
the exact finite formula
\begin{equation}
 R_\ell^{(y)}=
 \begin{cases}
 P_y(e_\ell^{(y)}-2e_{\ell+1}^{(y)}+e_{\ell+2}^{(y)}),&1\le\ell\le7,\\
 P_y e_8^{(y)},&\ell=8.
 \end{cases}
 \label{eq:run-formula}
\end{equation}
It also proves
\begin{equation}
 O_{y+1}=(s-1)O_y+L_y,\qquad s=p_{y+1}^2.
 \label{eq:odd-recurrence}
\end{equation}

\begin{lemma}[Per-run deletion ledger]
\label{lem:deletion}
Let an old positive run have length $1\le\ell\le8$, and let
$s=p_{y+1}^2$.  Across the $s$ replicated copies used to pass from
$w_y$ to $w_{y+1}$, the number of descendant even positive runs is
\begin{equation}
 \begin{cases}
 s-2,&\ell\text{ even},\\
 \ell-1,&\ell\text{ odd}.
 \end{cases}
 \label{eq:per-run-ledger}
\end{equation}
\end{lemma}

\begin{proof}
Represent each cyclic old run as an unwrapped interval bracketed by its two
old zero sites.  Those zeros remain zero in every lift, so no run crosses
the chosen $P_y$ seam and no descendant can merge across either bracket.
Because $(P_y,s)=1$, each old positive site has exactly one lift at which
the new prime square divides $2k+1$.  If two distinct positions $d,d'$ in
the same run were deleted in the same copy, subtraction would give
$2(d-d')\equiv0\pmod s$.  This is impossible because
$0<2|d-d'|\le14<s$ and every new $s=p_{y+1}^2$ is at least $25$.  Hence an
old run has $s-\ell$ untouched copies and exactly one single-site deletion
at each of its $\ell$ positions, all in distinct copies.

Suppose first that $\ell$ is even.  Each untouched copy contributes one
even run, giving $s-\ell$.  Deleting either endpoint leaves one odd run and
contributes no even run.  Each of the $\ell-2$ internal deletions leaves two
pieces of opposite parity, exactly one of which is a positive even run.
The total is $(s-\ell)+(\ell-2)=s-2$.

Now suppose that $\ell$ is odd.  Untouched copies contribute no even run.
For $\ell=1$, deleting the sole site leaves no positive piece, agreeing
with $\ell-1=0$.  Let $\ell\ge3$.  For a deletion at position
$d\in\{0,\ldots,\ell-1\}$, the two pieces have lengths $d$ and
$\ell-1-d$, hence the same parity.  Each of the two distinct endpoints
contributes one positive even piece, and every internal even deletion
position contributes two.  The total is
$2+2(\ell-3)/2=\ell-1$.  Odd deletion positions contribute none.  This
proves \eqref{eq:per-run-ledger}.
\end{proof}

\begin{proposition}[Even-run recurrence]
\label{prop:even-recurrence}
For every $y\ge1$,
\begin{equation}
 \boxed{\displaystyle
 \Ecal_{y+1}=(s-2)\Ecal_y+M_y,\qquad s=p_{y+1}^2.}
 \label{eq:even-recurrence}
\end{equation}
\end{proposition}

\begin{proof}
Sum Lemma~\ref{lem:deletion} over all old runs.  Each old even run supplies
$s-2$ descendant even runs.  An old odd run of length $\ell$ supplies
$\ell-1$, and the sum of those contributions is precisely $M_y$.
Old zero sites remain zero, while new deletions only split positive runs;
there is therefore no uncounted merger across old run boundaries.
\end{proof}

\section{Exact increment and strict monotonicity}

The locked RH-379 square-clock formula can be written
\begin{equation}
 G(q_y)=\frac{4}{\pi^2}
       +\frac{2O_y+4\Ecal_y}{A_y\pi^2}
       -\frac{\Ecal_y}{D_y}\kap.
 \label{eq:G-square}
\end{equation}
Put
\begin{equation}
 H_y=\kap\frac{A_y}{D_y}.
 \label{eq:H}
\end{equation}
RH-379 proves the exact Euler-product inequality
\begin{equation}
                         0<H_y<\frac4{\pi^2},\qquad
                         H_y\longrightarrow\frac4{\pi^2}.
 \label{eq:H-gap}
\end{equation}
No finite decimal is used in this statement.

\begin{theorem}[Exact square-clock increment]
\label{thm:increment}
For $s=p_{y+1}^2$,
\begin{equation}
\boxed{\begin{aligned}
G(q_{y+1})-G(q_y)
={}&\frac{2(L_y-2\Ecal_y)}{\pi^2A_y(s-1)}\\
&+\frac{M_y}{A_y(s-1)}
 \left(\frac4{\pi^2}-H_{y+1}\right).
\end{aligned}}
\label{eq:increment}
\end{equation}
Consequently
\begin{equation}
 G(q_{y+1})-G(q_y)
 \ge \frac{12}{\pi^2A_y(s-1)}>0.
 \label{eq:increment-lower}
\end{equation}
Thus $G(q_y)$ is strictly increasing for every $y\ge1$.
\end{theorem}

\begin{proof}
Besides \eqref{eq:odd-recurrence} and
\eqref{eq:even-recurrence}, the definitions give
\[
 A_{y+1}=(s-1)A_y,\qquad D_{y+1}=(s-2)D_y.
\]
Substituting all four recurrences into \eqref{eq:G-square} and subtracting
the $y$ expression gives, in the exact basis
$\mathbb Q\pi^{-2}+\mathbb Q\kap$,
\begin{equation}
 \frac{2(L_y-2\Ecal_y)+4M_y}{A_y(s-1)\pi^2}
 -\frac{M_y}{D_y(s-2)}\kap.
 \label{eq:increment-basis}
\end{equation}
Since
\[
 \frac{H_{y+1}}{A_y(s-1)}
 =\frac{\kap}{D_y(s-2)},
\]
equation \eqref{eq:increment-basis} is exactly
\eqref{eq:increment}.

The first numerator has the combinatorial identity
\begin{equation}
 L_y-2\Ecal_y
 =2R_4^{(y)}+4R_6^{(y)}+6R_8^{(y)}.
 \label{eq:X}
\end{equation}
RH-374 proves $R_8^{(y)}\ge1$, so this is at least six.  Also $M_y\ge0$
by definition, and the parenthesis in the second line of
\eqref{eq:increment} is positive by \eqref{eq:H-gap}.  Dropping that
nonnegative term yields \eqref{eq:increment-lower}.
\end{proof}

\begin{remark}[What is not monotone]
RH-379 writes $G(q_y)=B_y+\Delta_y$ and proves $\Delta_y\to0$.  Theorem
\ref{thm:increment} proves monotonicity of the sum $G(q_y)$; it neither uses
nor asserts monotonicity of $\Delta_y$.
\end{remark}

\section{A separator-specific same-support theorem}

The next result is not a general statement about covers of weighted cyclic
graphs.  Its proof uses the precise local densities and zero-weight phases
of the square clocks.

\begin{theorem}[Same-prime-support square-clock saturation]
\label{thm:saturation}
Fix $y\ge1$.  Suppose $q_y\mid Q$ and $Q$ has exactly the same prime
divisors as $q_y$.  Writing $Q=Rq_y$, one has
\begin{equation}
                         \boxed{G(Q)=G(q_y).}
 \label{eq:saturation}
\end{equation}
\end{theorem}

\begin{proof}
Every prime in either clock occurs to exponent at least two.  The exact
local squarefree formula therefore says that a phase has positive
$\dd$ precisely when it is nonzero modulo every supported prime square.
At such a phase,
\begin{equation}
 \dd_{q_y,r}=\frac{2}{A_y\pi^2},\qquad
 \dd_{Q,t}=\frac{2}{RA_y\pi^2}
 \quad(t\equiv r\pmod{q_y}),
 \label{eq:delta-scale}
\end{equation}
and both sides are zero simultaneously otherwise.  Likewise, the exact
two-squarefree formula gives
\begin{equation}
 \thh_{q_y,r}=\frac{\kap}{2D_y},\qquad
 \thh_{Q,t}=\frac{\kap}{2RD_y}
 \quad(t\equiv r\pmod{q_y}),
 \label{eq:theta-scale}
\end{equation}
whenever neither the current nor its predecessor is forced divisible by a
supported prime square, with simultaneous zero otherwise.  Thus every
three-state weight on a fine phase is exactly $1/R$ times the weight of its
projection modulo $q_y$.

Addition by two has an even and an odd cycle because both clocks are even.
On the even cycle, a phase divisible by $4$ has $\dd=\thh=0$.  On the odd
cycle, a phase divisible by $9$ has $\dd=\thh=0$.  At a zero-weight phase,
replace the action by $0$.  This cannot lower the score and, by
\eqref{eq:compatibility}, can only relax the two adjacent compatibility
conditions.  These mod-$4$ and mod-$9$ phases are therefore genuine
separators.

Under projection modulo $q_y$, each fine addition-by-two cycle traverses
the corresponding base cycle exactly $R$ times.  The separators cut it
into $R$ copies of the same finite path collection.  Each copied path has
weights scaled by $1/R$ by \eqref{eq:delta-scale}--\eqref{eq:theta-scale}.
Max-plus path optimization is positively homogeneous, so the $R$ copies
together have exactly the base optimum.  Summing the even and odd
components proves \eqref{eq:saturation}.  Finally, the RH-379 input
reflection transfers the positive optimum identity to the absolute value
defining $G$.
\end{proof}

\begin{remark}[Scoped negative control]
The prime-support hypothesis cannot be erased from this proof.  For
example, $36\mid180$, but $180$ adds the prime $5$, and the exact values are
\[
 G(36)=\frac{9}{2\pi^2}-\frac{\kap}{7},\qquad
 G(180)=\frac{73}{16\pi^2}-\frac{25\kap}{161}.
\]
The inequality is exact rather than a comparison of formal coefficient
pairs: with $H=\pi^2\kap<4$ from the locked RH-379 enclosure,
\[
 G(180)-G(36)
 =\frac1{\pi^2}\left(\frac1{16}-\frac{2H}{161}\right)
 >\frac1{\pi^2}\left(\frac1{16}-\frac8{161}\right)>0.
\]
This is a negative control against a general-multiple or general-cover
reading of Theorem~\ref{thm:saturation}; it is not a counterexample to the
theorem's same-support statement.
\end{remark}

\section{Finite-clock nonattainment and an explicit gap}

\begin{lemma}[Clock divisibility]
\label{lem:divisibility}
If $q\mid Q$, then $G(q)\le G(Q)$.
\end{lemma}

\begin{proof}
Repeat every phase table modulo $Q$ according to its residue modulo $q$.
The output word is unchanged for every input, so universal safety,
$c_{11}=0$, and the fixed-clock limiting score are unchanged.  The
$Q$-class can only have a larger optimum.  Input reflection again identifies
the positive and absolute formulations.
\end{proof}

\begin{theorem}[Quantitative finite-clock nonattainment]
\label{thm:nonattainment}
Let $q\ge1$ be any fixed finite clock.  Choose $y\ge1$ so that every odd
prime divisor of $q$ belongs to $\{p_1,\ldots,p_y\}$.  Then
\begin{equation}
 \boxed{\displaystyle
 G(q)<B_\infty,\qquad
 B_\infty-G(q)\ge
 \frac{12}{\pi^2A_y(p_{y+1}^2-1)}>0.}
 \label{eq:gap}
\end{equation}
Consequently no fixed finite clock attains the all-clock supremum
\eqref{eq:rh379-supremum}.
\end{theorem}

\begin{proof}
Set
\[
                         Q=\operatorname{lcm}(q,q_y).
\]
Then $q\mid Q$ and $q_y\mid Q$.  The prime support of $Q$ is exactly that
of $q_y$: the choice of $y$ contains all odd divisors of $q$, while the
factor $4$ in $q_y$ handles the prime $2$.  This remains true for an
arbitrary $2$-adic exponent in $q$ and arbitrary exponents of its supported
odd primes; those exponents merely enlarge $R=Q/q_y$.  Lemma
\ref{lem:divisibility} and Theorem~\ref{thm:saturation} give
\begin{equation}
                         G(q)\le G(Q)=G(q_y).
 \label{eq:lcm-chain}
\end{equation}

RH-379 proves $G(q_j)\to B_\infty$.  Theorem~\ref{thm:increment} makes
that sequence strictly increasing, hence $G(q_y)<B_\infty$.  More
quantitatively, telescope the nonnegative increments:
\[
 B_\infty-G(q_y)
 =\sum_{j\ge y}\bigl(G(q_{j+1})-G(q_j)\bigr)
 \ge G(q_{y+1})-G(q_y)
 \ge\frac{12}{\pi^2A_y(p_{y+1}^2-1)}.
\]
Combining this with \eqref{eq:lcm-chain} proves \eqref{eq:gap}.
\end{proof}

The quantifiers in Theorem~\ref{thm:nonattainment} are fixed-clock
quantifiers.  The index $y$ may be chosen after the finite integer $q$ is
given, but both $q$ and the auxiliary clock $Q$ are fixed before the
$N\to\infty$ limits already used to define $G$.  There is no diagonal
choice $q=q(N)$.

\section{Exact certificate and reproduction protocol}

The standard-library artifact is a proof companion, not a source of the
all-$y$ theorem.  It implements exact rational arithmetic in the basis
$u/\pi^2+v\kap$ and fails closed if the locked enclosure
$3.18<\pi^2\kap<3.19$ cannot decide a max-plus comparison.  The release
builder verifies that the sharper RH-379 rational interval lies strictly
inside this coarse interval.

The certificate performs four independent finite checks:
\begin{enumerate}[leftmargin=2em]
\item it compares the Euler-product run formula with direct cyclic words
      for $y=1,2,3$ and samples the per-run deletion ledger at
      $s=25,49,121$;
\item it checks the recurrence and increment algebra, including the exact
      anchors in Table~\ref{tab:anchors};
\item for nine same-support refinements it checks every fine-residue
      $\dd$ and $\thh$ scaling law, the mod-$4$/mod-$9$ separator causes,
      run replication, and an independent generic three-state cyclic
      max-plus dynamic program;
\item it locks $Q=180$ as a new-prime negative control and records finite
      least-common-multiple fixtures for the exponent scope.
\end{enumerate}

\begin{table}[ht]
\centering
\caption{Exact finite anchors.  These rows reproduce identities; the
all-order proofs are Theorem~\ref{thm:increment} and
Theorem~\ref{thm:saturation}.}
\label{tab:anchors}
\begin{tabular}{crrrrl}
\toprule
$y$&$\Ecal_y$&$L_y$&$M_y$&$L_y-2\Ecal_y$
 &$G(q_{y+1})-G(q_y)$\\
\midrule
1&1&8&0&6&$\frac{1}{16\pi^2}$\\[1mm]
2&23&160&24&114&$\frac{9}{256\pi^2}-\frac{24}{7567}\kap$\\[1mm]
3&1105&7160&1512&4950&$\frac{443}{30720\pi^2}-\frac{216}{128639}\kap$\\
\bottomrule
\end{tabular}
\end{table}

The result ledger has a recursively closed Draft 2020-12 schema.  Its 24
source locks exclude the mutable project instructions and handoff file.
For each locked predecessor input, the builder checks both the live SHA-256
and byte identity with the file at its declared release commit.  To replay
from this directory, run
\begin{verbatim}
make result
make test
make pdf
make archive
\end{verbatim}
The semantic PDF must be byte-identical to \texttt{main.pdf}; archive
verification fails on any missing member, hash mismatch, source-lock drift,
or PDF mismatch.

\section{Boundaries and next theorem edge}

The proved result is narrow but complete: it supplies an all-order even-run
recurrence, exact strict square-clock monotonicity, a
separator-specific saturation theorem, finite-clock nonattainment, and an
explicit positive gap within the frozen RH-379 class.

The following statements are outside the theorem.
\begin{itemize}[leftmargin=2em]
\item No table with phasewise $c_{11}(r)\ne0$ is covered.  The first missing
      object is a phase-weighted shift-two M\"obius correlation theorem.
\item Same-support saturation is not asserted for an arbitrary weighted
      cyclic graph, an arbitrary cover, or a multiplier that introduces a
      new prime.
\item No monotonicity is asserted for the correction $\Delta_y$.
\item No growing clock $q(N)$ and no convergence to adaptive distance-two
      capacity is asserted.
\item The construction supplies no canonical intrinsic dynamical spectral
      determinant, time-oriented scattering completion, self-adjoint
      generator, von Mangoldt prime-power trace identity, or completed-zeta
      divisor equality.  Gates A--E therefore remain false/open.
\item There is no Hilbert--P\'olya operator, no identification of Riemann
      zeros, and no implication for the Riemann hypothesis.
\end{itemize}

\paragraph{Immediate analytic rate trigger (not proved here).}
There is a within-class refinement beyond nonattainment.  Put
\[
 T_y=\sum_{p>p_y}\frac1{p^2-1},\qquad
 e_m=\prod_{p\ {\rm odd}}\left(1-\frac{m}{p^2}\right),
\]
and
\[
 X_\infty=
 \frac{2e_4-4e_5+6e_6-8e_7+10e_8}{e_1}>0.
\]
Indeed, if
$E_m^{(y)}=\prod_{i\le y}(1-m/p_i^2)$, the locked run formula gives
\[
 \frac{L_y-2\Ecal_y}{A_y}
 =\frac{2E_4^{(y)}-4E_5^{(y)}+6E_6^{(y)}
              -8E_7^{(y)}+10E_8^{(y)}}{E_1^{(y)}}
 \longrightarrow X_\infty.
\]
The same quantity is the normalized
$2R_4^{(y)}+4R_6^{(y)}+6R_8^{(y)}$, so
$X_\infty\ge6e_8/e_1>0$.  This identifies the proposed first-order
coefficient; it does not prove the tail asymptotic.
The exact reopen target is
\[
 B_\infty-G(q_y)=\frac{2X_\infty}{\pi^2}T_y+O(T_y^2),
\]
with the successor required to show that the $M_y$ contribution is only
second order.  RH-380 does not prove this rate theorem; proving the required
Euler-product expansion and $O(T_y^2)$ remainder, not computing further
finite rows or fitting decay, is the trigger.

\paragraph{First class-enlargement blocker.}
Extending beyond $c_{11}=0$ still requires a genuine theorem for
phase-weighted shift-two M\"obius correlations.  Without such an input, the
unrestricted-memory route remains \textsc{stop\_scoped}.

\section*{Declarations}

\paragraph{Data and code availability.}
All theorem inputs, exact code, tests, result schema, source locks, and
archive manifests are contained in the repository paper directory.  No
private dataset is used.

\paragraph{Author contributions.}
The RH research program performed conceptualization, formal analysis,
software construction, verification, visualization checks, and manuscript
preparation.

\paragraph{Funding.}
No external funding is declared.

\paragraph{Competing interests.}
No competing interests are declared.

\paragraph{Ethics and human participants.}
This mathematical and computational study uses no human participants,
animals, or personal data; ethics approval and consent are not applicable.

\paragraph{AI assistance disclosure.}
AI-assisted tools supported algebra checking, artifact generation, drafting,
and adversarial review.  Every released claim is constrained by the locked
repository sources and deterministic verification artifacts; responsibility
for the scoped mathematical statements remains with the research program.

\bibliographystyle{plain}
\bibliography{references}

\end{document}
